% Compact TikZ block diagram of the dual-corner architecture.
\begin{figure*}[t]
\centering
\resizebox{0.92\textwidth}{!}{%
\begin{tikzpicture}[
    node distance=1cm and 1.3cm,
    box/.style={rectangle, draw, rounded corners, minimum width=2.6cm,
                minimum height=0.8cm, align=center},
    group/.style={rectangle, draw, rounded corners, inner sep=5pt},
    arrow/.style={-{Latex[length=3mm]}, thick}
]

% Patient corner group (2x2 layout to save vertical space)
\node[group, label=above:{\textbf{Patient's Corner}}] (patientgroup) {
    \begin{tikzpicture}[node distance=0.4cm and 0.4cm]
        \node[box] (keypad) {$4\times8$ Alphanumeric\\Keypad};
        \node[box, right=of keypad] (temp) {LM35 Temp.\\Sensor};
        \node[box, below=of keypad] (pulse) {XD-58C\\Pulse Sensor};
        \node[box, right=of pulse] (plcd) {$16\times2$ LCD\\(4-bit parallel)};
    \end{tikzpicture}
};

% MCU
\node[box, right=3.4cm of patientgroup] (mcu)
     {ATmega32\\8~MHz internal RC\\(queue, ADC, UART)};

% Doctor corner group
\node[group, right=3.2cm of mcu, label=above:{\textbf{Doctor's Corner}}] (doctorgroup) {
    \begin{tikzpicture}[node distance=0.4cm]
        \node[box] (dlcd) {$16\times2$ LCD\\(I\textsuperscript{2}C)};
        \node[box, below=of dlcd] (btn) {Call-Next Button};
    \end{tikzpicture}
};

% Arrows
\draw[arrow] (patientgroup.east) --
      node[above, yshift=2pt, align=center]{Fields \&\\sensor samples}
      (mcu.west);
\draw[arrow] (mcu.east) --
      node[above, yshift=2pt, align=center]{Dequeued\\record}
      (doctorgroup.west);
\draw[arrow] (doctorgroup.south west) to[bend left=10]
      node[below, yshift=-2pt]{Dequeue} (mcu.south east);

% Optional wireless extension (not used by the delivered firmware)
\node[box, below=0.9cm of mcu] (bt) {Optional HC-05 / UART\\(future work)};
\draw[arrow, dashed] (mcu.south) -- (bt.north);

\end{tikzpicture}%
}
\caption{Dual-corner architecture. One ATmega32 mediates between the patient's
  corner (data entry and sensing) and the doctor's corner (review and queue
  advancement), so the queue is a single in-memory structure with one writer. The
  dashed UART path is wired out but unused.}
\Description{A block diagram with three groups arranged left to right. The left
group, labelled Patient's Corner, is a box containing four smaller boxes in a two
by two arrangement, holding a four by eight alphanumeric keypad, an LM35 temperature
sensor, an XD-58C pulse sensor, and a sixteen by two LCD in four-bit parallel
mode. A thick arrow labelled fields and sensor samples runs right from this group
to a central box labelled ATmega32, eight megahertz internal RC, queue, ADC and
UART. From the ATmega32 a thick arrow labelled dequeued record runs right to the
third group, labelled Doctor's Corner, which contains a sixteen by two I2C LCD
and a call-next button. A return arrow labelled dequeue curves back from the
Doctor's Corner to the ATmega32. Below the ATmega32, a dashed arrow points down
to a box labelled optional HC-05 slash UART, future work.}
\label{fig:system-architecture-block-diagram}
\end{figure*}


The system is organised as two interaction zones coordinated by one microcontroller
(Fig.~\ref{fig:system-architecture-block-diagram}). It is worth stating precisely what is
separated. The corners are distinct at the level of the user interface, since each has its own
display, input device and workflow, and neither can observe the other's intermediate state. They
are not separated at the level of computation, because a single ATmega32 scans the patient-side
keypad, samples both sensors, maintains the queue, polls the doctor's button and refreshes both
displays.

Consolidating onto one controller costs some independence and buys three things under R1 and R4. It
removes the need for an inter-controller link and its failure modes, it halves the controller bill of
materials, and, most importantly, it makes the queue a single in-memory structure with one writer. No
distributed-state problem remains to solve, since a record enqueued by the patient-side workflow is
immediately visible to the doctor-side workflow.

\subsection{Patient's Corner}
The patient interface is a self-service booth built from four peripherals. A $4\times8$ matrix
keypad, assembled from two commodity $4\times4$ membrane keypads, takes the name, age and mobile
number. An LM35 supplies temperature and an XD-58C optical sensor supplies pulse rate. A
$16\times2$ LCD in 4-bit parallel mode carries the prompts and the assigned serial number.

Registration is a strictly sequential dialogue, and the sequencing is what makes R3 achievable. The
LCD requests exactly one field at a time; the firmware will not advance until that field is accepted,
and will not accept a sensor reading until it has settled. A patient therefore never has to decide
what to do next, and cannot submit a record with a missing field. Only after every field is captured
does the firmware assign a serial number and enqueue the record, displaying the number last so the
patient leaves with a position they can reason about while waiting.

Placing the sensors in the booth rather than on the physician's desk is what recovers consultation
time. Temperature and pulse are acquired while the previous patient is still being seen, so
acquisition latency is hidden behind the consultation it would otherwise have extended, and
registration of later patients continues throughout.

\subsection{Doctor's Corner}
The physician's console is deliberately austere. It holds a $16\times2$ LCD on an
I\textsuperscript{2}C backpack, showing the current patient's serial number, name, age,
temperature and pulse, together with one momentary push button that calls the next patient.
There is no menu, no mode and no second control. A press either brings the patient at the head
of the queue onto the display or, once that patient has been seen, releases them and promotes
the next. If the queue is empty the display says so. The physician reads the serial number and calls
that patient in, already holding their vitals on screen. Reducing the console to one degree of freedom
answers R2 directly, leaving no navigation state to remember between patients and no way to reach an
unexpected screen mid-consultation.

Moving this display onto the I\textsuperscript{2}C bus is a necessity rather than a preference.
A second parallel LCD would consume six more I/O lines, and as Section~\ref{sec:hardware} shows,
the pin budget cannot absorb that once the 32-key matrix and two analog channels are accounted
for.

\begin{figure*}[t]
    \centering
    \includegraphics[width=\textwidth]{images/complete-hardware-architecture.png}
    \caption{Proteus schematic of the complete hardware architecture. The
      doctor-side I\textsuperscript{2}C LCD sits on PC0--PC1 and the
      patient-side 4-bit LCD on PD2--PD7, the LM35 and pulse sensor occupy
      ADC0 and ADC1, and the $4\times8$ keypad is driven from port B with row
      sensing on PA4--PA7.}
    \Description{A circuit schematic. At the left are a heart beat sensor module
    and an LM35 temperature sensor, whose analog outputs run to the ADC0 and ADC1
    pins of a central ATmega32 integrated circuit symbol. A pull-up resistor and
    a momentary push button connect to another port A pin. To the upper right of
    the microcontroller, a two-line I2C LCD module connects by just two wires
    labelled SCL and SDA; further right, a separate four-bit parallel LCD module
    connects by six wires from port D. Across the lower half of the figure is a
    grid of thirty-two push-button switches arranged in four rows and eight
    columns, labelled with digit-and-letter pairs such as 1/A and 2/B plus
    dedicated keys for CapsLock, NumLock, Space, BackSpace, Prev and Next. The
    eight column lines run to port B of the microcontroller and the four row
    lines run to port A through four ten-kilohm pull-down resistors.}
    \label{fig:proteus-complete}
\end{figure*}


Fig.~\ref{fig:proteus-complete} gives the schematic-level realisation, captured from the Proteus model
used for pre-deployment verification. It makes the shared-I/O structure concrete, showing the keypad
column drivers on port B, the row sense lines and both ADC inputs on port A, the
I\textsuperscript{2}C pair on port C, and the parallel LCD on port D.
